\documentclass{article}

% NeurIPS 2025 style -- user supplies neurips_2025.sty
\usepackage[final]{neurips_2024}

\usepackage[utf8]{inputenc}
\usepackage[T1]{fontenc}
\usepackage{hyperref}
\usepackage{url}
\usepackage{booktabs}
\usepackage{amsfonts}
\usepackage{amsmath}
\usepackage{nicefrac}
\usepackage{microtype}
\usepackage{xcolor}
\usepackage{array}
\usepackage{enumitem}

\title{Pre-Analysis Plan: Does Speculative Decoding Change What Benchmarks Measure?\\
\large Validity, Reliability, and Construct Invariance of LLM Evaluation Under Approximate Inference}

\author{%
  Rami Ratl Mrad\\
  Department of Computer Science\\
  Stanford University\\
  \texttt{ramimrad@stanford.edu}
}

\begin{document}

\maketitle

\begin{abstract}
Speculative decoding accelerates large language model (LLM) inference by having a small draft model propose token sequences that a larger target model then accepts or rejects. While theoretically lossless under exact acceptance sampling, practical deployments use approximate variants---relaxed thresholds, quantized draft models, batch-aware speculation---that may shift output distributions in subtle ways. We investigate whether these shifts are \emph{measurement-relevant}: do they affect benchmark scores, item-level response patterns, reliability estimates, or leaderboard rankings? Using a custom serving system, we conduct controlled comparisons between standard autoregressive decoding, exact speculative decoding, and approximate speculative decoding via probabilistic measurement models including 2PL IRT, Generalizability Theory, and ranking uncertainty analysis.
\end{abstract}

% ---------------------------------------------------------------
\section{Overview and Motivation}
% ---------------------------------------------------------------

\subsection{Problem Statement}

A growing fraction of LLM inference in both research and production is served through optimized inference engines that incorporate speculative decoding~\citep{leviathan2023fast, chen2023accelerating}. Practitioners and benchmark authors implicitly assume that inference method is a \emph{nuisance variable}---that switching from standard autoregressive decoding to speculative decoding leaves the underlying output distribution, and therefore benchmark scores, unchanged. This assumption is the foundation of score comparability across leaderboards, reproducibility of published results, and the validity of capability claims derived from benchmarks.

We argue this assumption deserves empirical scrutiny. Even when speculative decoding is theoretically distribution-preserving, practical implementations may introduce small but systematic differences through finite-precision arithmetic, batching behavior, relaxed acceptance criteria, or approximate draft verification. These effects may be negligible at the aggregate-score level while still altering item-level response patterns and benchmark reliability.

\textbf{Central research question:} Does speculative decoding, as implemented in practical serving systems, introduce statistically detectable and measurement-relevant differences in LLM benchmark performance. If so, do these differences vary by capability construct, draft model quality, and item difficulty?

\subsection{Connection to Course Topics}

This project directly engages the following CS321M topics:

\begin{itemize}[leftmargin=1.5em]
    \item \textbf{Item Response Theory (IRT):} We use 2PL IRT models and Differential Item Functioning (DIF) analysis to test whether item difficulty and discrimination parameters shift across decoding conditions.
    
    \item \textbf{Reliability and Generalizability Theory:} We decompose benchmark variance into item-, run-, and decoding-related components to quantify inference-induced measurement error.
    
    \item \textbf{Construct validity and invariance:} We evaluate whether benchmark-level capability relationships remain stable across decoding conditions.
    
    \item \textbf{Ranking uncertainty:} We analyze whether speculative decoding changes leaderboard rankings or increases rank uncertainty.
\end{itemize}

\subsection{Scope and Justification}

This is a one-person project. The author's custom LLM serving system enables controlled comparisons across decoding conditions, minimizing confounds from external frameworks. If speculative decoding introduces even small systematic biases, this could invalidate score comparisons across leaderboard entries using different inference backends---a direct implication for AI measurement in a world of diverse inference infrastructure.

% ---------------------------------------------------------------
\section{Literature Review}
% ---------------------------------------------------------------

\subsection{Speculative Decoding}

\citet{leviathan2023fast} introduced speculative decoding, showing that a small draft model can propose $K$ tokens simultaneously and a modified rejection sampling procedure recovers the exact target distribution. \citet{chen2023accelerating} independently developed a similar framework. Both prove theoretical losslessness under standard settings. Subsequent work has explored tree-structured speculation~\citep{miao2024specinfer}, $n$-gram caching, and relaxed acceptance criteria. Crucially, \emph{none examines benchmark measurement validity}---the assumption that output quality is preserved has not been tested at the item level using psychometric methods.

\subsection{Measurement Validity in NLP Evaluation}

\citet{ribeiro2020beyond} demonstrated that NLP benchmarks exhibit systematic failure modes invisible to aggregate accuracy. \citet{gururangan2018annotation} showed that many items contain spurious cues enabling models to succeed without genuine understanding---a construct validity problem. \citet{ethayarajh2022authenticity} applied IRT to NLP benchmarks, showing item difficulty parameters are stable and informative. Our work introduces a new measurement perturbation---the inference method itself---and asks whether IRT parameters are invariant across decoding conditions.

\subsection{Reliability and Reproducibility of LLM Evaluation}

\citet{dodge2020fine} documented substantial variance in LLM benchmark scores across random seeds, hyperparameters, and evaluation protocols. \citet{bouthillier2021accounting} extended this to a generalizability theory framework for machine learning evaluation. Most relevant to our work, \citet{liang2022holistic} (HELM) identified infrastructure-level sources of evaluation variance including sampling parameters and decoding strategies, but did not study speculative decoding specifically, as it was not yet widely deployed at the time of publication.

\subsection{Gap Addressed}

Prior work on speculative decoding focuses on inference efficiency and theoretical distribution preservation, while prior work on NLP measurement focuses on benchmark validity and evaluation instability. No existing work studies speculative decoding itself as a potential source of measurement error using psychometric methodology. Our project addresses this gap by evaluating whether inference method changes benchmark reliability, item behavior, or model rankings.

% ---------------------------------------------------------------
\section{Proposed Methods and Analysis}
% ---------------------------------------------------------------

\subsection{Experimental System}

The experimental platform is a basic LLM serving system adapted from a public repo which will feature:
\begin{itemize}[leftmargin=1.5em]
    \item Continuous batching
    \item PagedAttention for efficient KV-cache memory management
    \item Radix caching for prefix sharing across requests
\end{itemize}

This will be extended with speculative decoding supporting:
\begin{itemize}[leftmargin=1.5em]
    \item \textbf{Condition A (Baseline):} Standard autoregressive decoding, greedy or temperature-sampled
    \item \textbf{Condition B (Exact Speculation):} Draft model with rejection sampling recovering exact target distribution
    \item \textbf{Condition C (Approximate Speculation):} Relaxed acceptance threshold $\epsilon > 0$, trading distributional fidelity for throughput
\end{itemize}

Draft model size will be varied across experiments: we will use a family of models where a smaller variant (e.g., 1B parameter) acts as draft and a larger variant (e.g., 7B or 13B) acts as target. Acceptance rate $\alpha \in [0, 1]$ will be logged per request and used as a continuous covariate.

\subsection{Benchmark Selection}

We select benchmarks spanning four capability constructs to enable cross-construct comparisons:

\begin{center}
\begin{tabular}{llll}
\toprule
\textbf{Benchmark} & \textbf{Construct} & \textbf{Items} & \textbf{Format} \\
\midrule
MMLU~\citep{hendrycks2020measuring} & Factual knowledge & 14{,}042 & 4-way MCQ \\
GSM8K~\citep{cobbe2021training} & Mathematical reasoning & 1{,}319 & Free-form \\
HumanEval~\citep{chen2021evaluating} & Code generation & 164 & Function completion \\
HellaSwag~\citep{zellers2019hellaswag} & Commonsense reasoning & 10{,}042 & 4-way MCQ \\
\bottomrule
\end{tabular}
\end{center}

For each benchmark, we will collect item-level binary responses (correct/incorrect) for each model under each decoding condition across $R = 5$ independent runs (different random seeds for stochastic decoding; deterministic for greedy).

\subsection{Analysis Plan}

\paragraph{1. Aggregate Reliability Analysis.}
For each benchmark and decoding condition, we compute mean score with 95\% bootstrap CI, within-condition variance $\sigma^2_{\text{within}}$, and the between-condition difference $\Delta = \bar{s}_A - \bar{s}_B$ via paired $t$-test. The null hypothesis $\Delta = 0$ is tested per benchmark with Bonferroni correction.

\paragraph{2. Generalizability Theory Decomposition.}
We frame each benchmark evaluation as a $G$-study with facets: \emph{condition} (decoding method), \emph{item}, and \emph{run}. Using variance component estimation (restricted maximum likelihood), we partition total score variance into:
\begin{equation}
\sigma^2_{\text{total}} = \sigma^2_{\text{condition}} + \sigma^2_{\text{item}} + \sigma^2_{\text{run}} + \sigma^2_{\text{condition} \times \text{item}} + \sigma^2_{\text{error}}
\end{equation}
The $\sigma^2_{\text{condition}}$ and $\sigma^2_{\text{condition} \times \text{item}}$ terms quantify how much inference-method variance contaminates score interpretations. The $D$-study will characterize how many benchmark items are needed to achieve $G > 0.90$ under each condition.

\paragraph{3. Item Response Theory Analysis.}
We fit a 2-parameter logistic (2PL) IRT model to item-level response matrices under each decoding condition:
\begin{equation}
P(X_{ij} = 1 \mid \theta_j, a_i, b_i) = \frac{1}{1 + \exp(-a_i(\theta_j - b_i))}
\end{equation}
where $\theta_j$ is the latent ability of model $j$, $a_i$ is item discrimination, and $b_i$ is item difficulty. We use the \texttt{py-irt} library~\citep{lalor2019learning} for estimation.

We then conduct \textbf{Differential Item Functioning (DIF)} analysis treating decoding condition as the grouping variable. For each item, we test whether its IRT parameters ($a_i$, $b_i$) differ significantly between Condition A and Condition B using a likelihood ratio test with $\chi^2$ distribution. Items exhibiting significant DIF ($p < 0.05$ after FDR correction) are \emph{inference-sensitive items}: their measurement properties depend on how the model was run.

\paragraph{4. Construct-Level Stability Analysis.}

Rather than fitting item-level confirmatory factor models, which are impractical given the limited number of evaluated models, we analyze construct-level stability using benchmark-level correlations across decoding conditions. Specifically, we compare relative performance patterns across reasoning, factual recall, coding, and commonsense benchmarks under each inference condition.

We compute:
\begin{itemize}[leftmargin=1.5em]
    \item Correlations between benchmark scores across conditions
    \item Relative construct-specific score degradation
    \item Whether approximate speculation disproportionately affects reasoning-oriented benchmarks
\end{itemize}

This analysis provides a lightweight test of construct stability without requiring large-sample latent variable estimation.

\paragraph{5. Ranking Stability Analysis.}
We evaluate $M \geq 5$ models under all conditions, computing Kendall's $\tau$ between condition rankings, bootstrap CIs on each model's rank position, and identifying models whose rank CIs overlap under one condition but not the other (``ranking ambiguity'' induced by inference method).

\paragraph{6. Acceptance Rate as a Continuous Predictor.}
Using the approximate speculation condition (Condition C), we vary the acceptance threshold and measure the induced $\bar{\alpha}$ (mean acceptance rate per run). We fit a regression model:
\begin{equation}
\Delta s_{\text{benchmark}} = \beta_0 + \beta_1 (1 - \bar{\alpha}) + \epsilon
\end{equation}
to characterize how score degradation scales with deviation from the exact distribution. We test whether $\beta_1$ differs significantly across benchmark constructs (reasoning vs.\ factual recall), hypothesizing that multi-step reasoning tasks are more sensitive to approximate speculation.

\subsection{Hypotheses (Pre-registered)}

\begin{enumerate}[leftmargin=1.5em]
    \item \textbf{H1 (Aggregate Accuracy):} Under exact speculative decoding (Condition B), aggregate benchmark scores will not differ significantly from baseline (Condition A) at $\alpha = 0.05$.
    \item \textbf{H2 (Item-Level DIF):} A non-trivial fraction of items (estimated $> 5\%$) will exhibit DIF between conditions even when aggregate scores do not differ, indicating item-level distributional sensitivity invisible to summary statistics.
    \item \textbf{H3 (Construct-Specific Sensitivity):} Multi-step reasoning benchmarks (GSM8K) will show larger score degradation under approximate speculation (Condition C) than factual recall benchmarks (MMLU), due to compounding errors across token positions.
    \item \textbf{H4 (Measurement Invariance):} Metric invariance (equal factor loadings) will hold under exact speculation but fail under approximate speculation, indicating that the latent capability structure is sensitive to acceptance threshold.
    \item \textbf{H5 (Ranking Stability):} Model rankings will be stable across exact vs.\ baseline conditions but will show significant rank changes ($\tau < 0.90$) under approximate speculation with $\bar{\alpha} < 0.70$.
\end{enumerate}

\subsection{Evaluation Metrics}

\begin{itemize}[leftmargin=1.5em]
    \item Score-level: mean accuracy, 95\% bootstrap CI, paired $t$-test $p$-values, Cohen's $d$
    \item Reliability: Generalizability coefficient $G$, variance component estimates $\hat{\sigma}^2$
    \item IRT: Item fit statistics (RMSEA per item), DIF $\chi^2$ statistics, $\hat{\theta}_j$ estimate stability
    \item Invariance: CFI, RMSEA, $\Delta$CFI, $\Delta$RMSEA across invariance levels
    \item Ranking: Kendall's $\tau$, rank confidence interval width, rank inversion rate
\end{itemize}

% ---------------------------------------------------------------
\section{Timeline and Responsibilities}
% ---------------------------------------------------------------

This is a solo project. The timeline from the pre-analysis plan to the final manuscript is as follows:

\begin{center}
\begin{tabular}{p{2cm} p{9cm}}
\toprule
\textbf{Dates} & \textbf{Milestone} \\
\midrule
May 11--13 & Implement speculative decoding extension on existing serving system; validate exact-distribution recovery with a small synthetic test \\
May 13--16 & Run baseline (Condition A) benchmark evaluations across all four benchmarks; set up item-level logging \\
May 16--19 & Run Condition B and Condition C evaluations; verify acceptance rate logging; begin aggregate reliability analysis \\
May 19--22 & Fit 2PL IRT models; run DIF analysis; fit confirmatory factor models; run invariance tests \\
May 22--24 & Ranking stability analysis; acceptance rate regression; compile all results \\
May 24--27 & Write final manuscript; generate figures and tables \\
May 27 & \textbf{Final manuscript due} \\
May 27--June 3 & Clean and document code; write reproducibility README; submit \\
June 3 & \textbf{Code due} \\
\bottomrule
\end{tabular}
\end{center}

\paragraph{Identified Risks and Mitigations.}
\textbf{Compute constraints} limiting model count will be addressed by using small open-weight models (1B--7B) on a single GPU; $M \geq 5$ is achievable and sufficient for IRT.
\textbf{Implementation bugs} in speculative decoding will be caught by a KL-divergence validation test against the target distribution before any benchmark evaluation.
\textbf{IRT instability} from limited models will be mitigated by regularized estimation and treating DIF analyses as exploratory rather than confirmatory.
\textbf{Tight timeline} is managed by writing all analysis scripts before data collection begins.

% ---------------------------------------------------------------

\bibliographystyle{plainnat}
\bibliography{references}

\appendix

\section{Appendix: Draft Model Selection}

We will use the Qwen2.5 family (0.5B draft, 7B target) or Llama 3.2 (1B draft, 8B target), chosen because draft and target share a tokenizer (required for speculative decoding), both are openly licensed on HuggingFace, and the size ratio ($\approx$8--14$\times$) reflects typical production deployments.

\section{Appendix: Statistical Power Note}

The DIF analysis is primarily exploratory given the limited number of evaluated models. For MMLU (14,042 items) and HellaSwag (10,042 items), item-level patterns should be detectable; HumanEval (164 items) will be treated as secondary. We focus on breadth of item instability across conditions rather than precise latent ability estimation.

\end{document}